\chapter{Requirements}
\ChapterLabel{Requirements}
\index{requirement}
\index{requirement!implementation version}
\index{priority}

\section{Chapter Responsibility}
\label{Section:RequirementsResponsibility}
Specify the obligations the system must satisfy in a verifiable form.

\noindent\textbf{Expected result:} A controlled set of obligations that can be implemented and assessed.

\section{Revision 0.2 Scope Notice}
\label{Section:DeferredRequirements}
\begin{fsaNote}[Deferred project content --- template reference only]
This chapter retains template guidance and examples alongside explicitly
identified project requirement proposals and planning notes.
Its sample artifacts, IDs, priorities, statuses, versions, outcomes, and
approval examples do not describe the FTDI245R project. In particular, any
sample Done, Pass, or release entry is not implementation, verification,
validation, or approval evidence for this project. Detailed project content
is deferred under \hyperref[Section:ProjectScope]{the revision scope}.
See \hyperref[Table:ProjectGapReview]{the reviewed gap list} for actual
decisions and unresolved items.
\end{fsaNote}

\section{Draft Story Requirement Planning}
\label{Section:DraftStoryRequirements}
The original two proposals from \hyperref[Table:UserStories]{Section 4.5}
retain their stable IDs: \RequirementRef{REQ-D-002} and
\RequirementRef{REQ-F-004}. Joel Murphy (co-owner from Version 0.2: Jason Lefler) subsequently approved the acquisition,
timeout, disconnection, automatic-start/reconnection, and terminology decisions
recorded below, and authorized their translation into additional candidates.
The candidates are not an approved requirements baseline. Removed template
IDs remain reserved. \hyperref[Table:ProjectProfileConcernCoverage]{Profile
coverage} remains incomplete; unresolved decisions are owned by Joel Murphy and Jason Lefler.

\begingroup
\setlength{\emergencystretch}{2em}
\section{Project Requirements Summary}
\label{Section:ProjectRequirementsSummary}
\index{requirement!project proposals}
Project candidates are defined only in
\hyperref[Section:ProjectRequirements]{Section 5.14, Category-Specific Requirement Registers}.
Existing IDs are preserved; new candidates use the established Functional,
Quality, Interface, and Data families. Each candidate has proposed Must
priority, Open status, and \NoVersion\ for In Version. Joel Murphy and Jason Lefler own
priority/baseline review and all unresolved decisions. Tests are planned;
verification is  Pending. No implementation or product evidence
is claimed. Functional and quality subtypes are identified by the register
and obligation: read/control/recovery behavior, failure tolerance, raw-byte
interface, and data integrity.

A byte sequence is unchanged and in order when its length and every byte
position equal those of the comparison sequence. Delivery means acceptance
through the downstream consumer's interface, not completion of processing.
\Gls{fsareliable}, \gls{fsaplugandplay}, and \gls{fsacorruption} have
owner-approved definitions in the glossary. The completed-read predicate is
based on D2XX Section 3.10 \cite{FTDI_D2XX_v1_3}; every request is positive
under \RequirementRef{REQ-F-005}.

\paragraph{Agreed acquisition conditions and acceptance boundary.}
\phantomsection\label{Section:ProjectAcquisitionConditions}
\index{acquisition!agreed conditions}
Both software-generated known sequences and known sequences supplied through
physical FT245R hardware are required acceptance paths. Input is one unsigned
8-bit byte per ADC0804 conversion, with no headers or framing; the consumer
accepts the same raw bytes. Sampling is constant, with no source-rate bursts.
USB delivery may be batched; byte preservation spans batch boundaries.
Acquisition runs continuously until the user stops it. The sustained sample
rate is TBD pending ADC0804 timing and hardware confirmation; acceptance-test
duration, hardware/module revisions, supported Linux/Windows/macOS versions,
actual FTD2XX driver/library version, timeout duration, and quantitative
reliability targets are TBD (decision co-owners for each: Joel Murphy and Jason Lefler).
The D2XX Programmer's Guide v1.3 is the documentation baseline, not the
selected driver/library version. FTD2XX remains the proposed library.

On a supported host, installed driver/library and completed configuration
are prerequisites for automatic detection and acquisition on first connection.
Partial timeouts deliver only actual returned bytes in order and acquisition
continues; empty timeouts deliver no bytes and acquisition continues.
Disconnection does not crash the pipeline. All bytes successfully read into
host memory before disconnection remain subject to delivery and order
preservation. Bytes still in the device or generated while disconnected are
outside the preservation guarantee. The interruption is reported without
claiming complete acquisition; acquisition resumes automatically on
reconnection with the interruption clearly marked.

Consumer API/acceptance signaling, backpressure/failure handling, buffering
limits, device identity on reconnection, treatment of stale device-buffer
contents, reconnection timing, reporting mechanism, and user-stop/drain
details remain TBD (decision co-owners: Joel Murphy and Jason Lefler). These unresolved contracts
prevent a claim of established end-to-end feasibility or completed acceptance.
Known-sequence comparison covers the full source/consumer sequence in
uninterrupted tests and the successfully read host-memory sequence across
disconnect tests. Physical tests must retain source and read-boundary evidence;
software tests cannot substitute for target evidence.
Reciprocal links are maintained in
\hyperref[Table:ProjectRequirementTraceability]{Project Requirement Proposal Traceability}.


\paragraph{Installation and configuration prerequisites.}
\phantomsection\label{Section:ProjectStartupPrerequisites}
\index{acquisition!startup prerequisites}
For \RequirementRef{REQ-F-010}, prerequisites met means that the required
driver/library has been installed and configuration has been completed on a
supported host, as approved in
\hyperref[Section:ProjectAcquisitionConditions]{the acquisition conditions}.
FTD2XX remains the proposed library. Supported hardware/module revisions,
Linux/Windows/macOS versions, and actual driver/library version remain TBD
(decision co-owners: Joel Murphy and Jason Lefler). D2XX Programmer's Guide v1.3 is the documentation
baseline, not an installed library version.

The installation procedure, required configuration items and values, and
observable checks that establish successful installation and completed
configuration are TBD (decision co-owners: Joel Murphy and Jason Lefler). No specific settings,
installation checks, or completed deployment are implied by this precondition.

\paragraph{Retired template requirement mappings.}
\label{Section:RetiredTemplateRequirements}
The service-request template requirements have been removed from the active
registers. Links from retained example tests, traceability, and verification
records point to this notice instead of allocating those unrelated examples
to the project requirements. The removed template IDs are not reused.

\section{Embedded Profile Concern Coverage --- Open Completion Gate}
\label{Section:ProjectProfileConcernCoverage}
\index{domain tailoring profile!requirements coverage}
The \hyperref[Table:ProjectDomainProfiles]{Embedded Systems profile is Applicable}.
It is the only supplied profile; no additional Applicable or TBD profile was
identified. This register evaluates all 18 categories in
\hyperref[Table:EmbeddedRequirementCategories]{the profile checklist}, with
cross-cutting concerns below. Mapped means only the identified part is
addressed. Open entries are gaps, not Not Applicable dispositions, and do
not satisfy \hyperref[Section:DomainProfileRequirements]{the translation rule}.
All open decisions are owned by Joel Murphy and Jason Lefler. The requirements work is
\textbf{not complete}; the approved acquisition decisions do not waive the
remaining profile concerns.
\small
\begin{tabularx}{\linewidth}{L{3.4cm}L{4.2cm}Y}
\caption{Embedded Profile Concern Coverage\label{Table:ProjectProfileConcernCoverage}}\\
\toprule
\textbf{Profile concern} & \textbf{Requirement / disposition} & \textbf{Rationale and remaining decision}\\
\midrule
\endfirsthead
\multicolumn{3}{l}{\small\itshape Embedded Profile Concern Coverage (continued)}\\
\toprule
\textbf{Profile concern} & \textbf{Requirement / disposition} & \textbf{Rationale and remaining decision}\\
\midrule
\endhead
\bottomrule
\endlastfoot
Functional behavior & \RequirementRef{REQ-F-004}--\RequirementRef{REQ-F-013}; \RequirementRef{REQ-D-002} & Read recognition, positive requests, timeout continuation, continuous acquisition, detection/start, reconnection and interruption reporting are mapped. Other operational functions remain open.\\
Timing and real time & Partial: \RequirementRef{REQ-F-008} & Constant sampling; operation until user stop. Sample rate pending ADC0804 timing/hardware confirmation, test duration, timeouts, deadlines, latency and jitter remain TBD (decision co-owners: Joel Murphy and Jason Lefler).\\
Processing resources & Open & CPU, scheduling, execution-time, and host-processing budgets are TBD (decision co-owners: Joel Murphy and Jason Lefler).\\
Memory and storage & Open & Buffer capacity, overflow policy, memory/storage limits, and margins are TBD (decision co-owners: Joel Murphy and Jason Lefler). Sequence preservation is not a resource budget.\\
Power and thermal & Open & Supply/power domains, current, transitions, and thermal limits are TBD (decision co-owners: Joel Murphy and Jason Lefler).\\
Hardware interfaces & Partial: \RequirementRef{REQ-I-002}; \RequirementRef{REQ-D-002} & Unsigned 8-bit unframed raw bytes and delivery boundary defined. Electrical limits, pins, handshake ownership, units/scaling, controlled revisions and integration evidence remain TBD (decision co-owners: Joel Murphy and Jason Lefler).\\
Boot and initialization & Partial: \RequirementRef{REQ-F-009}; \RequirementRef{REQ-F-010} & Installed/configured prerequisites, automatic detection and first-connection start defined. Self-test, configuration validation, startup limits and failure response remain TBD (decision co-owners: Joel Murphy and Jason Lefler).\\
Reset and recovery & Partial: \RequirementRef{REQ-F-011}; \RequirementRef{REQ-D-002} & Automatic reconnection resume and preservation of host-read bytes defined. Retry/reset details, device identity, stale device bytes and recovery timing remain TBD (decision co-owners: Joel Murphy and Jason Lefler).\\
Safety & Open & Hazards, electrical protection, safe states, and approved safety dispositions are TBD (decision co-owners: Joel Murphy and Jason Lefler).\\
Security & Open & Threat/trust analysis, permissions, access control, physical/debug access, and data sensitivity are TBD (decision co-owners: Joel Murphy and Jason Lefler).\\
Reliability and availability & Partial: \RequirementRef{REQ-D-002}; \RequirementRef{REQ-Q-002} & Reliable/corruption definitions and disconnect preservation boundary approved. Device-resident and disconnected-source bytes excluded. Quantitative reliability/availability and resource feasibility remain TBD (decision co-owners: Joel Murphy and Jason Lefler).\\
Environmental & Open & Temperature, humidity, electromagnetic and other operating/storage conditions are TBD (decision co-owners: Joel Murphy and Jason Lefler).\\
Diagnostics and observability & Partial: \RequirementRef{REQ-F-004}; \RequirementRef{REQ-F-012}; \RequirementRef{REQ-F-013} & Read predicate and clearly marked interruption without complete-acquisition claim defined. Reporting mechanism, logs, counters and health interfaces remain TBD (decision co-owners: Joel Murphy and Jason Lefler).\\
Calibration and configuration & \hyperref[Section:ProjectSustainmentNotes]{Calibration records: approved Not Applicable} & Owner rationale: system not field deployed. Operational calibration needs, configurable values, integrity, and compatibility remain TBD (decision co-owners: Joel Murphy and Jason Lefler).\\
Firmware update & \hyperref[Section:ProjectSustainmentNotes]{Updates/rollback: approved Not Applicable} & Joel Murphy (co-owner from Version 0.2: Jason Lefler) disposition: system not field deployed. No new firmware implementation or update guarantee is inferred.\\
Manufacturing and provisioning & \hyperref[Section:ProjectSustainmentNotes]{Approved Not Applicable} & Joel Murphy (co-owner from Version 0.2: Jason Lefler) disposition: system not field deployed; manufacturing and field service excluded.\\
Regulatory and certification & Open & Applicable standards, evidence/independence obligations, and certification gates are TBD (decision co-owners: Joel Murphy and Jason Lefler).\\
Component lifecycle & \hyperref[Section:ProjectSustainmentNotes]{Component availability: approved Not Applicable} & Joel Murphy (co-owner from Version 0.2: Jason Lefler) disposition: system not field deployed. Exact supported parts/revisions and errata compatibility remain TBD (decision co-owners: Joel Murphy and Jason Lefler).\\
\end{tabularx}
\normalsize
\paragraph{Cross-cutting profile concerns.}
Hardware--software allocation, execution platform, interfaces, architecture
views, resource budgets, controlled toolchain, reproducible builds, host
compatibility versions, fixture calibration, target evidence, release
configuration, and evidence sufficiency remain
TBD (decision co-owners: Joel Murphy and Jason Lefler). The existing
\hyperref[Section:ProjectConfigurationControl]{Git/GitHub decision} is a
configuration-control fact, not proof of a reproducible implementation.
Field-service/update/manufacturing exclusions do not waive target-test,
safety, security, or operating-environment decisions. Every additional
applicable obligation needs a requirement mapping or a documented authorized
Not Applicable decision before the completion gate can close.

\endgroup

\section{How to Use This Chapter}
State singular, necessary, unambiguous, feasible, and verifiable obligations.
Every requirement shall have a stable ID, one priority, the version in which it
was implemented, at least one mapped use case, and at least one mapped test case.
Enter \NoVersion\ in \textbf{In Version} when no implementation version has
been assigned. Preserve IDs when wording or version assignments change.

Select the requirement categories that make the specification easier to review
and govern. Keep each requirement in one primary register and use traceability
links instead of copying the same requirement into multiple categories.

\section{Requirement Identifier Naming Convention}
\label{Section:RequirementIdentifierConvention}
\index{requirement!identifier convention}
\index{requirement!family code}
\index{requirement!scope qualifier}

Requirement identifiers use the following grammar:
\begin{verbatim}
REQ-[<scope>-]<family>-<sequence>
\end{verbatim}
The optional scope qualifier identifies a project-defined subsystem, product,
content type, or other governed area. Omit it for project-wide requirements.
The sequence is a three-digit, zero-padded number that is unique within the
exact scope and family combination. For example, \texttt{REQ-F-004} is the
fourth project-wide functional requirement, while
\texttt{REQ-API-Q-003} is the third quality requirement in a project-defined
\texttt{API} scope.

\small
\begin{tabularx}{\linewidth}{L{2.0cm}L{3.5cm}Y}
\caption{Template reference only: Requirement Identifier Segments and Family Codes\label{Table:RequirementIdentifierConvention}}\\
\toprule
\textbf{Code} & \textbf{Meaning} & \textbf{Use}\\
\midrule
\endfirsthead
\multicolumn{3}{l}{\small\itshape Requirement Identifier Segments and Family Codes (continued)}\\
\toprule
\textbf{Code} & \textbf{Meaning} & \textbf{Use}\\
\midrule
\endhead
\bottomrule
\endlastfoot
\texttt{REQ} & Requirement & Identifies the artifact as a requirement. Every requirement ID begins with this segment.\\
\texttt{<scope>} & Project-defined scope & Optional uppercase qualifier for a governed subsystem, product, content type, or other area. Define every scope code before use; absence denotes project-wide scope.\\
\texttt{B} & Business & States the stakeholder outcome, mission, commercial purpose, policy goal, or business result the system shall enable.\\
\texttt{F} & Functional & States observable behavior, service, workflow, calculation, state change, or capability the system shall provide.\\
\texttt{Q} & Quality & States a measurable quality attribute or non-functional expectation, including accessibility, performance, reliability, resilience, usability, maintainability, or compatibility.\\
\texttt{I} & Interface & States a user, software, API, hardware, communication, protocol, or external-system boundary contract, including integration behavior at that boundary.\\
\texttt{S} & Security & States security, privacy, consent, authorization, auditing, secret-handling, or protected-data obligations maintained in a specialized security register.\\
\texttt{D} & Data & States data definition, authority, provenance, quality, integrity, ownership, lifecycle, retention, or source-of-truth obligations.\\
\texttt{C} & Constraint & States a binding implementation, platform, technology, format, storage, source-control, standards, compatibility, schedule, budget, or governance constraint.\\
\texttt{\#\#\#} & Sequence & Three-digit, zero-padded stable sequence within the exact scope and family; assigned monotonically and never reused.\\
\bottomrule
\end{tabularx}
\normalsize

Classify a requirement by its primary obligation and verification concern. A
requirement may support several concerns, but its identifier contains only one
family code. Represent cross-cutting relationships through traceability rather
than a compound family code.

An AI agent shall reuse an existing scope and family whenever their definitions
fit. If no defined family accurately represents the requirement, the agent
shall not silently invent or use a new code. It shall propose a unique uppercase
code, full family name, decision boundary, rationale, and representative
examples; check the complete requirements register and category catalog for
semantic overlap and identifier collision; and obtain authorized-human
approval. After approval, the agent shall add the family to
Table~\ref{Table:RequirementIdentifierConvention}, update the category catalog
and AI guidance, and create or update the applicable register before assigning
an identifier from it. The same approval and documentation rule applies to a
new scope qualifier. Extending the convention shall not rename, renumber, or
reuse an existing identifier.

\section{Functional and Quality Requirements}
\index{requirement!functional}
\index{requirement!quality}
Functional and quality requirements answer different questions and shall not be
used interchangeably.

\begin{fsaNote}[Central Distinction]
A \textbf{functional requirement} states what behavior, service, calculation,
state transition, or externally observable result the system shall provide. A
\textbf{quality requirement} states how well that behavior or the system as a
whole shall perform under stated conditions. Quality requirements require a
measurable threshold, workload, environment, or other evaluation basis.
\end{fsaNote}

For example, ``The system shall display the current request status'' is
functional. ``The system shall display the current request status within two
seconds under the reference workload'' is a quality requirement. Avoid vague
quality statements such as ``the system shall be fast, secure, reliable, or
easy to use.'' Replace them with measurable conditions and success criteria.

Quality requirements are also called quality-attribute or non-functional
requirements. The term \emph{quality requirement} is preferred here because
these obligations describe behavior that shall be designed, implemented, and
verified rather than an absence of function.

\section{Requirement Status Model}
\index{requirement!status}
\index{status!Open}
\index{status!In Progress}
\index{status!Done}
\index{status!Retired}
Priority and status are independent. Priority records the requirement's
importance or commitment; status records its progress through the delivery
lifecycle. A status change shall not erase or replace the assigned priority.

\begin{tabularx}{\linewidth}{L{3.0cm}YL{3.3cm}}
\caption{Template reference only: Requirement Status Definitions\label{Table:RequirementStatusDefinitions}}\\
\toprule
\textbf{Status} & \textbf{Meaning} & \textbf{In Version Rule}\\
\midrule
\endfirsthead
\toprule
\textbf{Status} & \textbf{Meaning} & \textbf{In Version Rule}\\
\midrule
\endhead
\ReqStatusOpen & The requirement is active, but implementation has not started. Deferred work remains Open until work begins or the requirement is retired. & Enter \NoVersion.\\
\ReqStatusInProgress & Implementation or integration work has started, but the requirement does not yet satisfy the definition of Done. & Enter \NoVersion.\\
\ReqStatusDone & The requirement is implemented; mapped tests pass; verification is complete; traceability is current; and the completed behavior is included in a controlled version. & Enter the first controlled version containing the completed requirement.\\
\ReqStatusRetired & The requirement is no longer active because it was superseded, made obsolete, rejected, or withdrawn. Retain its stable ID, original text, and traceability; record the rationale and any replacement in document history. & Retain the original implementation version if it was previously Done; otherwise enter \NoVersion.\\
\bottomrule
\end{tabularx}

Only \ReqStatusDone\ asserts completion. Do not remove a retired requirement or
reuse its ID. A requirement may move from \ReqStatusRetired\ only through an
approved change that restores its active status and records the decision in
document history.

\section{Requirement Category Catalog}
\index{requirement!category catalog}
The following catalog distinguishes broad requirement families from their
subtypes. In particular, performance, capacity, scalability, reliability,
availability, resilience, recoverability, usability, accessibility,
maintainability, supportability, testability, portability, compatibility,
interoperability, security, privacy, safety, sustainability, and observability
are quality attributes. A requirement concerning one of these attributes is a
quality requirement when it specifies how well the system shall operate under
stated conditions.

A project may organize a specialized quality attribute, such as security,
privacy, or safety, in a separate register when it has distinct owners, review
criteria, evidence, or assurance obligations. The separate register is a
management view; it does not change the requirement's classification as a
quality requirement. Likewise, interface, data, operational, and other families
may contain both functional obligations and associated quality obligations. Use
the primary purpose of each requirement to select its register and use links to
represent cross-cutting concerns.

The catalog is intentionally broad. A project may use a family as a top-level
register, create registers for selected subtypes, combine closely related
groups, or omit groups that do not apply. Category choice shall not weaken the
traceability or verification rules for individual requirements.

\section{Domain Profile Impact on Requirements}
\label{Section:DomainProfileRequirements}
\index{requirement!domain profile}
Every Applicable domain profile shall be evaluated when developing or changing
requirements. Each applicable profile concern shall result in one or more
project requirements, a reference to an existing requirement, or a documented
Not Applicable decision with rationale approved by an authorized human. Record
the profile applicability decision in the Development Plan and retain
project-specific obligations only in the applicable requirement registers.
The profile itself is a tailoring guide, not a duplicate requirements
specification. See Section~\ref{Section:DomainProfiles} and the applicable
profile appendix.

\small
\begin{tabularx}{\linewidth}{L{3.0cm}L{4.0cm}Y}
\caption{Template reference only: Requirement Families, Subtypes, and Quality Attributes\label{Table:RequirementCategoryCatalog}}\\
\toprule
\textbf{Requirement Family} & \textbf{Subtype or Quality Attribute} & \textbf{Typical Scope}\\
\midrule
\endfirsthead
\multicolumn{3}{l}{\small\itshape Requirement Families, Subtypes, and Quality Attributes (continued)}\\
\toprule
\textbf{Requirement Family} & \textbf{Subtype or Quality Attribute} & \textbf{Typical Scope}\\
\midrule
\endhead
\midrule
\multicolumn{3}{r}{\small\itshape Continued on next page}\\
\endfoot
\bottomrule
\endlastfoot
Business & Outcomes and policy & Organizational outcomes, value, policy goals, measurable benefits, cost targets, and mission objectives.\\
Stakeholder and user & Stakeholder need & Capabilities, outcomes, and constraints needed by a stakeholder group or external party.\\
Stakeholder and user & User need & User goals and externally visible outcomes expressed independently of a particular implementation.\\
Functional & System behavior & Services, calculations, decisions, state changes, workflows, and observable system responses.\\
Quality & General quality & Measurable characteristics that state how well a function or the system shall operate under stated conditions.\\
Quality & Performance & Response time, latency, throughput, processing time, and resource consumption.\\
Quality & Capacity & Supported data volume, transaction volume, concurrent users, devices, or retained records.\\
Quality & Scalability & Behavior as workload, data, users, or deployed instances increase or decrease.\\
Quality & Reliability & Probability or frequency of failure-free operation under stated conditions.\\
Quality & Availability & Required service availability, planned downtime, and permitted outage duration.\\
Quality & Resilience and fault tolerance & Continued or degraded operation when components, dependencies, or communication paths fail.\\
Quality & Recoverability and continuity & Recovery time, recovery point, restoration, backup, failover, and disaster-continuity objectives.\\
Quality & Usability & Task effectiveness, efficiency, learnability, error prevention, and user satisfaction.\\
Quality & Accessibility & Conformance and measurable support for users with visual, auditory, motor, speech, or cognitive disabilities.\\
Quality & Maintainability & Time, effort, modularity, analyzability, and change characteristics required for system modification.\\
Quality & Supportability & Diagnostic access, service procedures, field replacement, support hours, and support-team capabilities.\\
Quality & Testability & Controllability, observability, repeatability, test interfaces, and evidence needed to verify behavior.\\
Quality & Portability & Ability to install, move, or operate the system across specified environments or platforms.\\
Quality & Compatibility & Coexistence and correct operation with specified versions, devices, browsers, operating systems, or products.\\
Quality & Interoperability & Coordinated behavior, protocols, semantics, orchestration, and exchanges with external systems.\\
Quality & Security & Authentication, authorization, confidentiality, integrity, cryptography, auditing, and threat resistance.\\
Quality & Privacy & Collection, consent, purpose, minimization, access, disclosure, retention, deletion, and rights for personal data.\\
Quality & Safety & Hazard prevention, risk reduction, safe state, interlocks, alarms, and unacceptable-loss controls.\\
Quality & Sustainability & Energy use, resource efficiency, carbon reporting, hardware life, repair, reuse, and disposal objectives.\\
Quality & Observability & Logs, metrics, traces, health signals, diagnostic context, and measurable diagnostic coverage.\\
Interface & User interface & Human-system interaction, presentation, navigation, input, feedback, and interaction contracts.\\
Interface & Software and API & Operations, messages, schemas, errors, protocols, sequencing, and software boundary contracts.\\
Interface & Hardware and communication & Connectors, signals, electrical characteristics, device protocols, networks, and physical boundaries.\\
Data and information & Data definition and quality & Data definitions, schemas, provenance, accuracy, completeness, integrity, ownership, and permitted transformations.\\
Data and information & Lifecycle and governance & Creation, access, retention, archival, legal hold, disposition, deletion, and proof of destruction.\\
Operational & Operation and administration & Startup, shutdown, scheduling, configuration, provisioning, monitoring, backup, recovery, and incident response.\\
Operational & Deployment and platform & Approved platforms, topology, packaging, installation, configuration, upgrade, and runtime environment.\\
Operational & Infrastructure and environment & Network, compute, storage, facility, power, temperature, physical location, and environmental limits.\\
Transition & Migration and rollout & Data conversion, coexistence, rollout, cutover, rollback, legacy integration, and decommissioning.\\
Compliance and legal & Regulatory and standards & Obligations imposed by regulations, standards, certifications, policies, or assurance frameworks.\\
Compliance and legal & Legal and contractual & Licenses, warranties, intellectual-property terms, service agreements, procurement terms, and contractual duties.\\
Compliance and legal & Reporting and audit & Reports, evidence, audit trails, review access, content, timing, and record accountability.\\
Product support & Documentation & Required user, operator, administrator, developer, API, maintenance, and compliance documentation.\\
Product support & Training & Required curricula, exercises, qualifications, delivery modes, and demonstrated learner outcomes.\\
Commercial & Licensing and entitlement & Product editions, entitlement, metering, third-party components, subscriptions, and licensing restrictions.\\
AI and machine learning & Model and data controls & Model performance, data provenance, bias evaluation, explainability, human oversight, drift, and retraining controls.\\
Physical and hardware & Physical characteristics & Dimensions, mass, materials, connectors, tolerances, environmental ratings, and hardware behavior.\\
Constraint & Solution constraint & Mandated technologies, designs, budgets, schedules, standards, dependencies, or solution boundaries.\\
\end{tabularx}
\normalsize

\section{Creating a Project-Specific Category}
Create an additional category when it gives requirements a distinct meaning,
owner, review method, or assurance process. Define the category's scope and a
unique family code using
Section~\ref{Section:RequirementIdentifierConvention}; duplicate one of the
category-specific registers below; keep the standard columns; add the category to
\hyperref[Table:RequirementCategoryCatalog]{the category catalog}; and update
use cases, tests, traceability, verification, validation, document history, and
the index as applicable. Document whether the new category is top-level or a
subtype of an existing category.

Do not create a category solely to avoid resolving an ambiguous requirement.
If one statement contains multiple independently verifiable obligations, split
it into singular requirements and assign each one a primary category.

\section{Review Rules}
Reject a requirement row if its priority, status, use-case mapping, or test-case
mapping is empty. The \textbf{Status} field shall contain Open, In Progress,
Done, or Retired. The \textbf{In Version} field shall follow
\hyperref[Table:RequirementStatusDefinitions]{the status definitions}. Confirm
that reciprocal mappings appear in the user-story, use-case, test,
traceability, and verification registers.

\begin{landscape}
\small
\section{Category-Specific Requirement Registers}
\label{Section:ProjectRequirements}
The seven category registers retain their existing column format and
reference labels. Functional, Quality, Interface, and Data registers contain
the project candidates. Previous template rows have been removed.
The other registers remain pending profile review, not Not Applicable;
decision co-owners: Joel Murphy and Jason Lefler. Each requirement is defined in one register only.

\subsection{Business Requirements}
\index{requirement!business}
\begin{tabularx}{\linewidth}{L{2.1cm}YL{2.0cm}L{2.1cm}L{1.5cm}L{2.6cm}L{2.6cm}}
\caption{Business Requirements Register\label{Table:BusinessRequirements}}\\
\toprule
\textbf{ID} & \textbf{Requirement} & \textbf{Priority} & \textbf{Status} & \textbf{In Version} & \textbf{Use Case(s)} & \textbf{Test Case(s)}\\
\midrule
\endfirsthead
\multicolumn{7}{l}{\small\itshape Business Requirements Register (continued)}\\
\toprule
\textbf{ID} & \textbf{Requirement} & \textbf{Priority} & \textbf{Status} & \textbf{In Version} & \textbf{Use Case(s)} & \textbf{Test Case(s)}\\
\midrule
\endhead
\midrule
\multicolumn{7}{r}{\small\itshape Continued on next page}\\
\endfoot
\bottomrule
\endlastfoot
\multicolumn{7}{l}{\textit{No project requirement proposed yet; profile review pending. Decision co-owners: Joel Murphy and Jason Lefler.}}\\
\end{tabularx}

\clearpage
\subsection{Functional Requirements}
\begin{tabularx}{\linewidth}{L{2.1cm}YL{2.0cm}L{2.1cm}L{1.5cm}L{2.6cm}L{2.6cm}}
\caption{Functional Requirements Register\label{Table:FunctionalRequirements}}\\
\toprule
\textbf{ID} & \textbf{Requirement} & \textbf{Priority} & \textbf{Status} & \textbf{In Version} & \textbf{Use Case(s)} & \textbf{Test Case(s)}\\
\midrule
\endfirsthead
\multicolumn{7}{l}{\small\itshape Functional Requirements Register (continued)}\\
\toprule
\textbf{ID} & \textbf{Requirement} & \textbf{Priority} & \textbf{Status} & \textbf{In Version} & \textbf{Use Case(s)} & \textbf{Test Case(s)}\\
\midrule
\endhead
\midrule
\multicolumn{7}{r}{\small\itshape Continued on next page}\\
\endfoot
\bottomrule
\endlastfoot
\RequirementID{REQ-F-004} & The pipeline shall recognize a read as completed if and only if its returned status is \texttt{FT\_OK} and its actual returned byte count equals its requested byte count. & \gls{must} & \ReqStatusOpen & \NoVersion & \UseCaseRef{UC-002} & \TestCaseRef{TC-011}\\
\RequirementID{REQ-F-005} & The pipeline shall request at least one byte in every acquisition read. & \gls{must} & \ReqStatusOpen & \NoVersion & \UseCaseRef{UC-002} & \TestCaseRef{TC-012}\\
\RequirementID{REQ-F-006} & The pipeline shall continue acquisition after a partial read timeout returning \texttt{FT\_OK}. & \gls{must} & \ReqStatusOpen & \NoVersion & \UseCaseRef{UC-002} & \TestCaseRef{TC-013}\\
\RequirementID{REQ-F-007} & The pipeline shall continue acquisition after an empty read timeout returning \texttt{FT\_OK}. & \gls{must} & \ReqStatusOpen & \NoVersion & \UseCaseRef{UC-002} & \TestCaseRef{TC-013}\\
\RequirementID{REQ-F-008} & During uninterrupted operation under the agreed acquisition conditions, the pipeline shall continue acquisition until the user stops it. & \gls{must} & \ReqStatusOpen & \NoVersion & \UseCaseRef{UC-001} & \TestCaseRef{TC-014}\\
\RequirementID{REQ-F-009} & On a supported host with the required driver/library installed and configuration completed, the pipeline shall automatically detect connected FT245R hardware without a code change. & \gls{must} & \ReqStatusOpen & \NoVersion & \UseCaseRef{UC-001}; \UseCaseRef{UC-003} & \TestCaseRef{TC-014}\\
\RequirementID{REQ-F-010} & The pipeline shall automatically start acquisition on first connection of the supported FT245R hardware after the \hyperref[Section:ProjectStartupPrerequisites]{installation and configuration prerequisites} are met. & \gls{must} & \ReqStatusOpen & \NoVersion & \UseCaseRef{UC-001}; \UseCaseRef{UC-003} & \TestCaseRef{TC-014}\\
\RequirementID{REQ-F-011} & Following disconnection, the pipeline shall automatically resume acquisition when the supported FT245R hardware reconnects. & \gls{must} & \ReqStatusOpen & \NoVersion & \UseCaseRef{UC-002}; \UseCaseRef{UC-004} & \TestCaseRef{TC-015}\\
\RequirementID{REQ-F-012} & The pipeline shall clearly report an acquisition interruption caused by FT245R disconnection. & \gls{must} & \ReqStatusOpen & \NoVersion & \UseCaseRef{UC-002}; \UseCaseRef{UC-004} & \TestCaseRef{TC-015}\\
\RequirementID{REQ-F-013} & The pipeline shall not represent acquisition spanning a disconnection as complete and uninterrupted. & \gls{must} & \ReqStatusOpen & \NoVersion & \UseCaseRef{UC-002}; \UseCaseRef{UC-004} & \TestCaseRef{TC-015}\\
\end{tabularx}

\subsection{Quality Requirements}
\begin{tabularx}{\linewidth}{L{2.1cm}YL{2.0cm}L{2.1cm}L{1.5cm}L{2.6cm}L{2.6cm}}
\caption{Quality Requirements Register\label{Table:QualityRequirements}}\\
\toprule
\textbf{ID} & \textbf{Requirement} & \textbf{Priority} & \textbf{Status} & \textbf{In Version} & \textbf{Use Case(s)} & \textbf{Test Case(s)}\\
\midrule
\endfirsthead
\multicolumn{7}{l}{\small\itshape Quality Requirements Register (continued)}\\
\toprule
\textbf{ID} & \textbf{Requirement} & \textbf{Priority} & \textbf{Status} & \textbf{In Version} & \textbf{Use Case(s)} & \textbf{Test Case(s)}\\
\midrule
\endhead
\midrule
\multicolumn{7}{r}{\small\itshape Continued on next page}\\
\endfoot
\bottomrule
\endlastfoot
\RequirementID{REQ-Q-002} & The pipeline shall remain running when the FT245R disconnects. & \gls{must} & \ReqStatusOpen & \NoVersion & \UseCaseRef{UC-002}; \UseCaseRef{UC-004} & \TestCaseRef{TC-015}\\
\end{tabularx}

\clearpage
\subsection{Interface Requirements}
\index{requirement!interface}
\begin{tabularx}{\linewidth}{L{2.1cm}YL{2.0cm}L{2.1cm}L{1.5cm}L{2.6cm}L{2.6cm}}
\caption{Interface Requirements Register\label{Table:InterfaceRequirements}}\\
\toprule
\textbf{ID} & \textbf{Requirement} & \textbf{Priority} & \textbf{Status} & \textbf{In Version} & \textbf{Use Case(s)} & \textbf{Test Case(s)}\\
\midrule
\endfirsthead
\multicolumn{7}{l}{\small\itshape Interface Requirements Register (continued)}\\
\toprule
\textbf{ID} & \textbf{Requirement} & \textbf{Priority} & \textbf{Status} & \textbf{In Version} & \textbf{Use Case(s)} & \textbf{Test Case(s)}\\
\midrule
\endhead
\midrule
\multicolumn{7}{r}{\small\itshape Continued on next page}\\
\endfoot
\bottomrule
\endlastfoot
\RequirementID{REQ-I-002} & The pipeline shall accept an unframed stream containing one unsigned 8-bit byte per ADC0804 conversion, with no headers. & \gls{must} & \ReqStatusOpen & \NoVersion & \UseCaseRef{UC-001} & \TestCaseRef{TC-016}\\
\end{tabularx}

\subsection{Security Requirements}
\index{requirement!security}
\begin{tabularx}{\linewidth}{L{2.1cm}YL{2.0cm}L{2.1cm}L{1.5cm}L{2.6cm}L{2.6cm}}
\caption{Security Requirements Register\label{Table:SecurityRequirements}}\\
\toprule
\textbf{ID} & \textbf{Requirement} & \textbf{Priority} & \textbf{Status} & \textbf{In Version} & \textbf{Use Case(s)} & \textbf{Test Case(s)}\\
\midrule
\endfirsthead
\multicolumn{7}{l}{\small\itshape Security Requirements Register (continued)}\\
\toprule
\textbf{ID} & \textbf{Requirement} & \textbf{Priority} & \textbf{Status} & \textbf{In Version} & \textbf{Use Case(s)} & \textbf{Test Case(s)}\\
\midrule
\endhead
\midrule
\multicolumn{7}{r}{\small\itshape Continued on next page}\\
\endfoot
\bottomrule
\endlastfoot
\multicolumn{7}{l}{\textit{No project requirement proposed yet; profile review pending. Decision co-owners: Joel Murphy and Jason Lefler.}}\\
\end{tabularx}

\clearpage
\subsection{Data Requirements}
\index{requirement!data}
\begin{tabularx}{\linewidth}{L{2.1cm}YL{2.0cm}L{2.1cm}L{1.5cm}L{2.6cm}L{2.6cm}}
\caption{Data Requirements Register\label{Table:DataRequirements}}\\
\toprule
\textbf{ID} & \textbf{Requirement} & \textbf{Priority} & \textbf{Status} & \textbf{In Version} & \textbf{Use Case(s)} & \textbf{Test Case(s)}\\
\midrule
\endfirsthead
\multicolumn{7}{l}{\small\itshape Data Requirements Register (continued)}\\
\toprule
\textbf{ID} & \textbf{Requirement} & \textbf{Priority} & \textbf{Status} & \textbf{In Version} & \textbf{Use Case(s)} & \textbf{Test Case(s)}\\
\midrule
\endhead
\midrule
\multicolumn{7}{r}{\small\itshape Continued on next page}\\
\endfoot
\bottomrule
\endlastfoot
\RequirementID{REQ-D-002} & Under the \hyperref[Section:ProjectAcquisitionConditions]{agreed acquisition conditions}, the pipeline shall deliver every byte successfully read into host memory to the downstream processing consumer unchanged and in its original read order, including bytes pending delivery at disconnection. & \gls{must} & \ReqStatusOpen & \NoVersion & \UseCaseRef{UC-001}; \UseCaseRef{UC-002}; \UseCaseRef{UC-004} & \TestCaseRef{TC-010}; \TestCaseRef{TC-013}; \TestCaseRef{TC-015}; \TestCaseRef{TC-016}\\
\end{tabularx}

\subsection{Constraint Requirements}
\index{requirement!constraint}
\begin{tabularx}{\linewidth}{L{2.1cm}YL{2.0cm}L{2.1cm}L{1.5cm}L{2.6cm}L{2.6cm}}
\caption{Constraint Requirements Register\label{Table:ConstraintRequirements}}\\
\toprule
\textbf{ID} & \textbf{Requirement} & \textbf{Priority} & \textbf{Status} & \textbf{In Version} & \textbf{Use Case(s)} & \textbf{Test Case(s)}\\
\midrule
\endfirsthead
\multicolumn{7}{l}{\small\itshape Constraint Requirements Register (continued)}\\
\toprule
\textbf{ID} & \textbf{Requirement} & \textbf{Priority} & \textbf{Status} & \textbf{In Version} & \textbf{Use Case(s)} & \textbf{Test Case(s)}\\
\midrule
\endhead
\midrule
\multicolumn{7}{r}{\small\itshape Continued on next page}\\
\endfoot
\bottomrule
\endlastfoot
\multicolumn{7}{l}{\textit{No project requirement proposed yet; profile review pending. Decision co-owners: Joel Murphy and Jason Lefler.}}\\
\end{tabularx}
\end{landscape}

\index{acquisition!continuous operation}
\index{read timeout!partial and empty}
\index{disconnection!preservation boundary}
\index{reconnection!automatic resume}
\index{plug and play}
\index{corruption!byte preservation}
